\bookchapter{第三十五章：兒童與手機——數字時代嘅教育指南}

\chapterintro{本章唔講「手機係洪水猛獸」，亦唔鼓吹「越早接觸越聰明」。我哋按年齡分段畀設備、設規則、教素養、護健康，一步步嚟。}

\booksection{一、幾歲畀仔手機？}

冇標準答案，但有年齡基線。過早畀手機，仔大腦前額葉尚未發育成熟，自控力跟不上算法嘅「誘惑設計」；過晚畀，又可能等佢喺同伴中「社交脫節」。下面係分齡建議：

\begin{itemize}
  \item \textbf{3-6歲}：盡量唔碰手機。如需使用，單次唔超過15分鐘，內容以家長精選嘅教育動畫為主，且必須\textbf{陪同觀睇}。呢一階段螢幕對視力與注意力發育影響最大。
  \item \textbf{7-9歲}：可喺家長設備上「借用」，用於網課、聽故事、查資料。開始建立「使用前申請、用完歸還」嘅規則，唔配置專屬手機。
  \item \textbf{10-12歲}：可考慮畀一臺「功能機」或兒童手錶，解決上下學聯絡。如確需智能機，應開啟嚴格嘅家長控制。
  \item \textbf{13歲以上}：進入初中，社交需求增強，可以畀智能機，但需配合「使用公約」同定期復盤。
\end{itemize}

\begin{tcbinfo}{「夠用就好」原則}
仔嘅第一部手機唔需要旗艦。處理器嘅「過剩效能」往往被用嚟跑遊戲。選一臺電池夠大、效能中等、螢幕護眼嘅千元機或官方翻新機，反而更有利管控。預算參考：1500-2500元區間即可滿足絕大多數學生需求。

\end{tcbinfo}

\booksubsection{智能手錶作為過渡方案}

如果淨係想解決「放學後能不能聯絡上」嘅問題，兒童手錶比手機更合適。佢淨係可以打電話、發語音、定位，冇瀏覽器同短影片，天然限制咗「沉迷入口」：

\begin{itemize}
  \item \textbf{小天才電話手錶}：生態最完整，同齡人社交「碰一碰加好友」係其核心賣點，但專屬App生態較封閉。
  \item \textbf{華為兒童手錶}：同華為手機聯動好，定位精準，支援視像通話，性價比高。
  \item \textbf{小米兒童手錶}：價格親民，基礎功能齊全，適合預算有限嘅家庭。
\end{itemize}

買手錶時注意：選擇支援\textbf{4G全網通}嘅型號（2G/3G正逐步退網），並確認運營商兒童套餐資費。

\booksection{二、家長控制系統全攻略}

設備畀到仔之前，第一件事唔係開機，而係設定家長控制。兩大平臺都有原生方案：

\booksubsection{iOS：螢幕使用時間}

\begin{itemize}
  \item \textbf{App限額}：設定 > 螢幕使用時間 > App限額，可為每個App類別設定每日可用時長（如遊戲30分鐘）。
  \item \textbf{停用時間}：設定每日「熄燈時段」（如22:00-7:00），期間所有App變灰禁用，淨係電話可用。
  \item \textbf{內容限制}：可按年齡段過濾影視、網頁、App安裝，並禁止App內購買。
  \item \textbf{家庭共享}：用家長嘅iPhone遠程管理仔嘅設備，批准下載請求、查看使用報告。
\end{itemize}

\booksubsection{安卓：Google Family Link / 品牌家長控制}

\begin{itemize}
  \item \textbf{Google Family Link}：谷歌官方方案，可遠程管理設備使用時長、鎖定設備、查看活動報告。國內安卓需視機型可用性。
  \item \textbf{小米兒童空間}：設定 > 兒童空間，可設定使用時長、白名單App，退出需家長密碼。
  \item \textbf{華為學生模式}：設定 > 健康使用手機 > 學生模式，限制使用時段與App，支援遠程管理。
  \item \textbf{OPPO/vivo兒童模式}：喺「設定」或專屬「兒童空間」App中開啟，原理類似。
\end{itemize}

\begin{center}
\begin{tabularx}{\linewidth}{|X|X|X|}
\hline
\textbf{功能} & \textbf{iOS 螢幕使用時間} & \textbf{安卓 Family Link / 品牌方案} \tabularnewline
\hline
App限額 & 支援，按類別 & 支援，按App \tabularnewline
停用時間 & 支援 & 支援（品牌方案略有差異） \tabularnewline
內容過濾 & 系統級，較徹底 & 依賴品牌，參差不齊 \tabularnewline
遠程管理 & 家庭共享 & Family Link / 品牌雲 \tabularnewline
安裝App審批 & 需家長批准 & Family Link 支援 \tabularnewline
\hline
\end{tabularx}
\end{center}

\begin{tcbwarn}{家長控制唔係「一勞永逸」}
仔總能搵到繞過嘅方法（網上搜「點樣破解螢幕使用時間」嘅教程一大把）。技術手段淨係底線，真正嘅防線係親子溝通同信任。建議每月同仔一齊復盤使用報告，而唔係單方面監控。

\end{tcbwarn}

\booksection{三、青少年模式}

各大內容平臺都內置咗「青少年模式」，開啟後會限制使用時長、過濾唔良內容、關閉充值打賞。呢個係國家法規要求嘅強制功能，但\textbf{需要家長主動開啟並設定密碼}：

\begin{itemize}
  \item \textbf{抖音/快手青少年模式}：每日限40分鐘（22:00-6:00禁用），內容淨係展示精選科普、教育類影片，無法直播、無法打賞、無法充值。
  \item \textbf{微信青少年模式}：限制微信支付（需家長授權額度），限制公眾號、影片號內容，搖一搖、附近嘅人等功能預設關閉。
  \item \textbf{遊戲防沉迷系統}：國家層面嘅實名認證+時間限制。未成年人僅可喺週五、週六、週日及法定節假日嘅20:00-21:00登入遊戲，每日1小時。
\end{itemize}

\textbf{實名認證唔好用家長身份證「借號」}\par

好多仔用家長身份證註冊遊戲帳號繞過防沉迷，呢個既違反用戶協議，亦等仔失去保護。家長應保管好自己嘅支付密碼與身份資訊，唔好為咗慳事「借號」畀仔。

\booksection{四、網絡素養教育}

技術能攔截90\%嘅風險，剩低10\%靠仔嘅判斷力。比起「唔畀上網」，更重要嘅係教佢哋「點樣上網」：

\begin{itemize}
  \item \textbf{私隱保護}：唔向陌生人透露真實姓名、學校、家庭住址；唔發送個人相片、身份證、家庭合影；定位功能預設關閉，僅喺使用時開啟。
  \item \textbf{網絡欺凌識別與應對}：被辱罵、孤立、惡意P圖都屬於欺凌。話畀仔聽遇到時唔好對罵、唔好刪除證據、第一時間話畀家長聽。微信可長按消息投訴，平臺會處理施暴帳號。
  \item \textbf{內容批判}：教仔分辨「事實」同「觀點」，識別標題黨、震驚體；理解算法淨係推你想睇嘅，唔代表世界嘅全貌；遇到「天上跌餡餅」嘅廣告一律唔信。
  \item \textbf{網絡安全}：唔撳陌生連結、唔掃來歷不明嘅二維碼、唔下載非官方商店嘅App；密碼唔分享畀任何人（包括「好朋友」）；遇到索要驗證碼嘅一律視為騙子。
\end{itemize}

\begin{tcbpractice}{動手：同仔一齊簽一份《手機使用公約》}
將規則寫低、雙方簽字貼喺顯眼處，比口頭約定有效得多。建議包含：

\begin{itemize}
  \item 使用時段：如「作業完成後19:30-20:30可用」。
  \item 禁用區域：餐桌上、臥室睡前、課堂上手機一律收起。
  \item 內容邊界：唔安裝未告知家長嘅App，唔同陌生人私聊。
  \item 違約後果：如「超時一次，次日使用時長減半」。
  \item 家長承諾：唔喺仔面前「掃手機」超過約定時長，以身作則。
\end{itemize}

\end{tcbpractice}

\booksection{五、健康使用習慣}

仔嘅眼睛、頸椎、睡眠都喺發育期，手機對身體嘅傷害比成年人更大：

\begin{itemize}
  \item \textbf{螢幕時間建議}：3-6歲每日唔超過1小時；6-12歲學習類唔限制，娛樂類唔超過1小時；12歲以上同仔協商，但娛樂類唔超過2小時。
  \item \textbf{護眼習慣}：遵守「20-20-20法則」——每睇螢幕20分鐘，抬頭睇20英尺（約6米）外嘅遠處20秒；睡前1小時唔睇手機，藍光會抑制褪黑素影響睡眠。
  \item \textbf{姿勢管理}：手機螢幕與眼睛平齊，唔好低頭含胸；連續使用唔超過30分鐘就要起身活動。
  \item \textbf{親子共同規則}：餐桌、臥室、書房設「無手機區」；全家一齊睇戲、戶外活動時手機統一收起，等「放低手機」成為家庭文化而唔係對仔嘅懲罰。
\end{itemize}

\begin{tcbinfo}{家長係最好嘅「防火牆」}
研究表明，父母自身手機使用時長與仔呈強正相關。要求仔唔睇手機，自己卻掃短影片到深夜，規則就會失效。改變仔，先由放低自己嘅手機開始。

\end{tcbinfo}

\booksection{六、數字遺產與兒童}

呢一節容易被忽略，但好重要。仔嘅數字資產（帳號、相片、充值餘額、遊戲裝備）同樣面臨繼承與私隱問題：

\begin{itemize}
  \item \textbf{兒童數字資產嘅繼承}：仔不幸離世時，其社交帳號、遊戲帳號、雲盤相片嘅繼承流程與成年人相同，需提供死亡證明、親屬關係證明。建議家長平時就為仔嘅重要帳號做「遺產聯絡人」或「緊急聯絡人」設定（如蘋果遺產聯絡人、微信帳號繼承）。
  \item \textbf{未成年人嘅數字私隱}：《個人資訊保護法》規定，處理不滿14周歲未成年人個人資訊需取得監護人同意，且應制定專門嘅私隱規則。家長應定期審查仔常用App嘅私隱政策，關閉「個性化推薦」同「數據共享」開關。
  \item \textbf{帳號註冊建議}：為仔註冊嘅帳號應使用家長可控嘅郵箱/手機號，便於喺未成年階段統一管理；仔成年後可再做帳號移交。
\end{itemize}

\begin{tcbwarn}{警惕「兒童數字遺產」嘅另一面}
仔由出世起就喺度被「記錄」——家長曬嘅成長相、幼兒園發嘅影片、智能手錶嘅健康數據。呢啲數據一旦上傳就難以徹底刪除。發布前諗一諗：十年後仔會唔會因為呢張相感到尷尬？尊重仔嘅數字人格，由剋制「曬娃」開始。

\end{tcbwarn}

\begin{tcbsummary}{本章小結}
\begin{itemize}
  \item 畀手機嘅年齡冇標準答案，但3-6歲盡量唔碰、13歲前用兒童手錶或功能機過渡
  \item 「夠用就好」：第一部手機選千元機而唔係旗艦，效能過剩等於誘惑過剩
  \item iOS螢幕使用時間與安卓Family Link/品牌家長控制係兩套技術防線，配合使用更穩
  \item 青少年模式與遊戲防沉迷係法規要求，但需家長主動開啟並守護密碼
  \item 網絡素養比技術攔截更重要：教仔識別欺凌、保護私隱、批判內容
  \item 健康習慣要由家長做起——你放低手機，仔先至會放低
  \item 兒童數字遺產與私隱係新興議題，平時就應做「遺產聯絡人」設定並剋制曬娃
\end{itemize}

\end{tcbsummary}
